\section{Formal Model}
\label{sec:semantics}

This section gives an operational semantics for \hadl that alternates evaluation, type checking, and policy checking, with self-healing when any check fails. Types, workflows, and authorization policies materialize at runtime and reduction deliberately visits ill-typed intermediate states during self-heal, so we (i) parameterize over a non-deterministic oracle, (ii) give a substitution-based de-Bruijn core so that materialized values flow by $\beta$-reduction rather than through a typed environment, and (iii) state soundness relative to an oracle that is eventually-truthful within the declared retry budget.

\subsection{Machine Configuration and Judgments}
\label{sec:sem-config}

A \hadl configuration is a four-tuple $C = \MConfigS{\MErr}{\MPol}{\MStore}{e}$ where $\MErr$ is a \emph{heal context} holding only failed-attempt feedback events $\ErrEv{\epsilon}$ --- flushed to $\emptyset$ on a successful oracle step and extended by $\cdot \ErrEv{\epsilon}$ on a heal step --- whose retry metric $\retries{\MErr}$ is simply $|\MErr|$; $\MPol$ is the installed Cedar policy, updated only by $\Enforce{v}$ via an append-only merge that only tightens; $\MStore$ is a finite map from mutable-variable names to $(\tau, v)$ pairs, mutated only by \texttt{var} declarations and assignments; and $e$ is the expression under reduction. The store $\MStore$ is orthogonal to the substitution-based de-Bruijn core: \texttt{let}, closure application, and $\rulename{For-Cons}$ use capture-avoiding substitution, while \texttt{var} declarations, assignments, and reads use $\MStore$; the two surfaces do not interact. Rules that leave $\MStore$ unchanged elide it and use the three-component notation $\MConfig{\MErr}{\MPol}{e}$; rules that read or write it are collected in \sect{app:rules-mut}. The current principal lives inside the $\Gen{\tau}{s}{\MPrin}$, $\Agent{s}{\MPrin}$, and closure-application syntax. Expressions use de-Bruijn indices and carry their own binders; values are the subset satisfying $\mathsf{isValue}$ (base literals, records, arrays, schema values $\Tschema(\tau)$, Cedar policy values, closures).

Four judgments drive correctness. \emph{Runtime typing} $\RtType{v}{\tau}$ holds of closed values (substitution keeps every checked value closed). \emph{Residual static typing} is black-boxed: it is re-invoked at runtime after every materialization exactly as at compile time; only the reflexivity cases the soundness proof uses are exposed. \emph{Policy authorization} $\PolOk{\MPol}{\MPrin}{a}$ lifts the Cedar decision procedure for $a \in \{\texttt{gen},\texttt{agent}\}$. \emph{Well-formedness} $\WF{C}$ requires both $\retries{\MErr} \le \MRetry$ and that every store cell is well-typed: $\forall x,\tau,v.\; \MStore(x) = (\tau, v) \Rightarrow \RtType{v}{\tau}$. Residual typeability and environment well-typedness now live inside the reduction rules as $\RtType{\cdot}{\cdot}$ side conditions that substitution composes.

\subsection{Reduction}
\label{sec:sem-reduction}

The reduction relation $C \step C'$ is a single small-step relation parameterized by the oracle $\mathcal{O}$, split between \emph{root-level} rules that fire when the top of $e$ matches and \emph{structural congruence} rules that drive reduction to the leftmost sub-term under CBV. Let-binding is $\beta$-reduction at de-Bruijn index $0$: $\rulename{Let-Bind}$ rewrites $\Let{\tau}{v}{e}$ to $\Inst{e}{v}$, the capture-avoiding substitution replacing $\mathrm{var}\;0$ by $v$; closure application $\rulename{Eval-Success}$ performs the analogous simultaneous substitution.

Self-healing is \emph{let-centric} and \emph{continuation-driven}. The expression $\Gen{\tau}{s}{\MPrin}$ is no longer a standalone redex: it can only reduce as the immediate right-hand side of a $\texttt{let}$. We classify materialization targets by the predicate $\Healable{\tau}$ --- true on Schema, Policy, arrow types, and any record/array type containing a healable component; false on the base scalars Unit, Bool, Number, String. At each healable $\tau$ the runtime re-checks the oracle's response \emph{and} re-typechecks the continuation $p$ under the extended context $\Gext{x}{\tau}$; the second of those checks is what drives self-heal. We write both as the single-valued judgment $\Gctx \vdash t : \tau \Rightarrow r$ with $r \in \{\mathsf{ok}, \mathsf{err}(\epsilon)\}$: it applies to values (the runtime re-check) and to expressions (the continuation check). In the error branch the judgment \emph{returns} the feedback string $\epsilon$ consumed by self-healing. The context $\Gctx$ resolves schema-name references bound by type declarations (e.g.\ validating $v$ against a previously-installed \texttt{user\_schema}); it is threaded outward unchanged from the enclosing program and we elide it when no confusion arises. Each healable $\tau$ has exactly two rules: $\rulename{Let-Gen-Success}_\tau$ fires when $v$ is well-typed at $\tau$ \emph{and} the continuation $p$ is well-typed at its declared return type $\tau'$ under $\Gext{x}{\tau}$, stepping to $\Inst{p}{v}$ and flushing $\MErr$ to $\emptyset$; $\rulename{Let-Gen-Heal}_\tau$ fires when $v$ is well-typed but the continuation is not, extending $\MErr$ with the continuation's $\epsilon$ and retrying the same redex. If the oracle returns an ill-typed $v$, no rule fires --- the configuration is stuck by omission. Two uniform rules apply at \emph{every} $\tau$: $\rulename{Let-Gen-TypeError}$ steps an ill-typed response at non-healable $\tau$ to a fresh sink value $\errV$ that is intentionally untypeable ($\TJok{\Gctx}{\errV}{\tau}$ is unprovable for every $\tau$, fail-closed), and $\rulename{Let-Gen-HealPol}$ records a feedback event and retries when the policy denies the action. \figr{fig:rules-core} shows the Schema pair and \figr{fig:rules-workflow} the arrow (workflow) pair, both following the same continuation-driven pattern; for arrows, $\Inst{p}{v}$ is higher-order substitution of the oracle-returned closure into the continuation, mechanized via the \texttt{lean-subst} framework. The full rule set --- including the $\rulename{Let-Gen-BudgetError}$ sink and the Policy pair, which is subsumed by the parametric $\rulename{Let-Gen-Success-Healable}$/$\rulename{Let-Gen-Heal-Healable}$ rules via $\Healable{\Tpolicy} = \mathsf{true}$ --- appears in Appendix~\ref{app:rules}.

Because $\errV$ has no runtime type, T2 (staged materialization, \sect{thm:hadl-soundness}) is stated as ``$v \neq \errV \Rightarrow \exists \tau.\ \RtType{v}{\tau}$''. The predicate $\RtType{v}{\tau}$ is shorthand for the $\mathsf{ok}$-branch $\TJok{\emptyset}{v}{\tau}$ on closed values and is used wherever the feedback string is irrelevant (e.g., T2 and $\rulename{Let-Bind}$). The standalone $\Agent{s}{\MPrin}$ keeps its own pair $\rulename{Agent-Success}$/$\rulename{Agent-Heal-Pol}$: agent's materialization type $\Tstring$ is non-healable, so no type-heal rule is needed; conceptually $\Agent{s}{\MPrin}$ is shorthand for $\Let{\Tstring}{\Agent{s}{\MPrin}}{\unit}$.

\begin{figure}[t]
\small
\begin{mathpar}
\inferrule*[right=\rulename{Let-Gen-Success-Schema}]
  {\PolOk{\MPol}{\MPrin}{\texttt{gen}} \\
   \Oracle{s}{\MErr}{\Tschema} \ni v \\
   \TJok{\Gctx}{v}{\Tschema} \\
   \TJok{\Gext{x}{\Tschema}}{p}{\tau'}}
  {\MConfig{\MErr}{\MPol}{\Let{\Tschema}{\Gen{\Tschema}{s}{\MPrin}}{p}}
   \step
   \MConfig{\emptyset}{\MPol}{\Inst{p}{v}}}

\inferrule*[right=\rulename{Let-Gen-Heal-Schema}]
  {\PolOk{\MPol}{\MPrin}{\texttt{gen}} \\
   \Oracle{s}{\MErr}{\Tschema} \ni v \\
   \TJok{\Gctx}{v}{\Tschema} \\
   \TJerr{\Gext{x}{\Tschema}}{p}{\tau'}{\epsilon} \\
   \retries{\MErr \cdot \ErrEv{\epsilon}} \le \MRetry}
  {\MConfig{\MErr}{\MPol}{\Let{\Tschema}{\Gen{\Tschema}{s}{\MPrin}}{p}}
   \step
   \MConfig{\MErr \cdot \ErrEv{\epsilon}}{\MPol}{\Let{\Tschema}{\Gen{\Tschema}{s}{\MPrin}}{p}}}
\end{mathpar}
\caption{Self-healing for the Schema materialization target. Both rules require a well-typed oracle response ($\TJok{\Gctx}{v}{\Tschema}$); they differ in whether the \emph{continuation} $p$ type-checks under $\Gext{x}{\Tschema}$. $\rulename{Let-Gen-Success-Schema}$ flushes the heal context to $\emptyset$ and commits; $\rulename{Let-Gen-Heal-Schema}$ appends the continuation's feedback $\epsilon$ and retries the same redex. If the oracle returns an ill-typed $v$, no rule fires --- the configuration is stuck by omission, and the progress lemma (T4c) discharges this case under an eventually-truthful oracle. The arrow (workflow) pair is shown in \figr{fig:rules-workflow}; the parametric pair $\rulename{Let-Gen-Success-Healable}$/$\rulename{Let-Gen-Heal-Healable}$, given in the appendix, generalizes both pairs to any $\Healable{\tau}$ and additionally covers nested compound healable types such as $\Tarr{\Tschema}$ (clinical-trial \texttt{Array[crf]}) and $\Tpolicy$ (whose segment-authorization check is the existing $\PolOk{}{}{\texttt{gen}}$ premise of the parametric rule).}
\label{fig:rules-core}
\end{figure}

\begin{figure}[t]
\small
\begin{mathpar}
\inferrule*[right=\rulename{Let-Gen-Success-Arrow}]
  {\PolOk{\MPol}{\MPrin}{\texttt{gen}} \\
   \Oracle{s}{\MErr}{\Tarrow{T_1,\ldots,T_n}{T'}} \ni v \\
   \TJok{\Gctx}{v}{\Tarrow{T_1,\ldots,T_n}{T'}} \\
   \TJok{\Gext{x}{\Tarrow{T_1,\ldots,T_n}{T'}}}{p}{\tau''}}
  {\MConfig{\MErr}{\MPol}{\Let{\Tarrow{T_1,\ldots,T_n}{T'}}{\Gen{\Tarrow{T_1,\ldots,T_n}{T'}}{s}{\MPrin}}{p}}
   \step
   \MConfig{\emptyset}{\MPol}{\Inst{p}{v}}}

\inferrule*[right=\rulename{Let-Gen-Heal-Arrow}]
  {\PolOk{\MPol}{\MPrin}{\texttt{gen}} \\
   \Oracle{s}{\MErr}{\Tarrow{T_1,\ldots,T_n}{T'}} \ni v \\
   \TJok{\Gctx}{v}{\Tarrow{T_1,\ldots,T_n}{T'}} \\
   \TJerr{\Gext{x}{\Tarrow{T_1,\ldots,T_n}{T'}}}{p}{\tau''}{\epsilon} \\
   \retries{\MErr \cdot \ErrEv{\epsilon}} \le \MRetry}
  {\MConfig{\MErr}{\MPol}{\Let{\Tarrow{T_1,\ldots,T_n}{T'}}{\Gen{\Tarrow{T_1,\ldots,T_n}{T'}}{s}{\MPrin}}{p}}
   \step
   \MConfig{\MErr \cdot \ErrEv{\epsilon}}{\MPol}{\Let{\Tarrow{T_1,\ldots,T_n}{T'}}{\Gen{\Tarrow{T_1,\ldots,T_n}{T'}}{s}{\MPrin}}{p}}}
\end{mathpar}
\caption{Self-healing for the arrow (workflow) materialization target. The pair is structurally identical to \figr{fig:rules-core} with $\tau = \Tarrow{T_1,\ldots,T_n}{T'}$: both rules require a well-typed oracle response and diverge on whether the continuation $p$ type-checks under $\Gext{x}{\Tarrow{T_1,\ldots,T_n}{T'}}$. On success the substitution $\Inst{p}{v}$ commits the generated closure to the continuation --- higher-order substitution of a workflow value into a program, mechanized as a straightforward instance of the \texttt{lean-subst} framework already underpinning $\rulename{Let-Bind}$ and closure application. On continuation failure the feedback $\epsilon$ returned by the type checker is appended to $\MErr$ and the same redex is retried. An ill-typed $v$ is stuck by omission, discharged by progress (T4c) under an eventually-truthful oracle.}
\label{fig:rules-workflow}
\end{figure}

The pure-core rules ($\rulename{If-True/False}$, $\rulename{While}$, $\rulename{For-Nil/Cons}$, $\rulename{Js}$, $\rulename{Seq}$, $\rulename{Ask}$, $\rulename{Say}$, $\rulename{Enforce-Install}$) form a standard CBV relation; $\rulename{For-Cons}$ uses $\beta$-reduction to bind the iteration variable. $\rulename{Enforce-Install}$ acts on a policy value at the head; the install operation rejects blanket $\texttt{forbid}$ under a deny-all guard, rejects unknown action types, and merges rules append-only. Allow-set shrinkage ($\Allow{\MPol'} \subseteq \Allow{\MPol}$) is an algebraic consequence of install, mechanized as $\mathtt{policyInstall\_shrinks}$.

\subsection{Correctness}
\label{sec:sem-correctness}

Classical progress and preservation do not apply directly: types are runtime values, policies evolve, and heal steps deliberately visit ill-typed configurations. We decompose what \emph{is} preserved into three per-step \emph{Soundness} components plus \emph{Progress} for oracle redexes, all parameterized by the oracle, and argue \emph{Termination} on paper.

\begin{definition}[Eventually-truthful oracle]
\label{def:oracle-truthful}
An oracle $\mathcal{O}$ is \emph{eventually-truthful at budget $\MRetry$} for a site $(s,\tau)$ under an authorization predicate $A : V \to \mathrm{Prop}$ if there exists a heal context $\MErr$ with $\retries{\MErr} \le \MRetry$ and a value $v$ satisfying $\Oracle{s}{\MErr}{\tau} \ni v$, $\RtType{v}{\tau}$, and $A(v)$. We instantiate $A(v) \equiv \top$ in the soundness theorems below; $A$ lets callers (e.g., agent sites) rule out self-contradictory responses without weakening the progress lemma.
\end{definition}

This matches the existential shape of the Lean predicate $\mathtt{Oracle.eventuallyTruthful}$: it witnesses some well-typed response at \emph{some} heal context, not the caller's current $\MErr$.

\begin{theorem}[Soundness]
\label{thm:hadl-soundness}
For any oracle $\mathcal{O}$ and any step $C \step C'$:
(\textbf{T1}, WF-preservation) $\WF{C} \Rightarrow \WF{C'}$; (\textbf{T2}, staged materialization) if $C' = \MConfigS{\MErr'}{\MPol'}{\MStore'}{v}$ with $\mathsf{isValue}(v)$ and $v \neq \errV$, then $\exists \tau.\ \RtType{v}{\tau}$; (\textbf{T3}, policy monotonicity) $\Allow{\MPol'} \subseteq \Allow{\MPol}$.
\end{theorem}

T2 is conditional on $v \neq \errV$ because $\errV$ is the fail-closed sink produced by $\rulename{Let-Gen-TypeError}$ and $\rulename{Let-Gen-BudgetError}$ and has no runtime type by design. For non-sink values T2 forces the mechanization to assign runtime types to closures, records, and arrays via dedicated $\mathtt{vClos}/\mathtt{vRec}/\mathtt{vArr}$ constructors that black-box structural content; the $\RtType{v}{\Tschema}$ premise of $\rulename{Let-Gen-Success-Schema}$ still witnesses the annotated type at oracle materialization.

\begin{theorem}[Progress for oracle redexes]
\label{thm:hadl-progress}
Let $C = \MConfigS{\MErr}{\MPol}{\MStore}{e}$ with $\WF{C}$. \textbf{T4a} (no heal past budget for standalone gen): if $e = \Gen{\tau}{s}{\MPrin}$ and $\retries{\MErr} > \MRetry$, then $C$ has no successor at all (standalone $\Gen{\tau}{s}{\MPrin}$ is not a redex). \textbf{T4b} (truthful success at Schema): if $e = \Let{\Tschema}{\Gen{\Tschema}{s}{\MPrin}}{\mathrm{var}\;0}$, $\PolOk{\MPol}{\MPrin}{\texttt{gen}}$, and $\mathcal{O}$ is eventually-truthful at $(s,\Tschema)$, then $\exists C'.\ C \step C'$ via $\rulename{Let-Gen-Success-Schema}$. \textbf{T4c} (progress-gen-Schema, progress-agent): with the same authorization hypothesis and a direct oracle witness at the caller's $\MErr$, $\exists C'.\ C \step C'$.
\end{theorem}

T4a is now \emph{stronger} than ``budget-exhausted gen is stuck'': because $\Gen{\tau}{s}{\MPrin}$ is no longer a standalone redex, no rule fires from such an expression at any heal context (regardless of budget). Materialization happens only inside a $\Let{\tau}{\Gen{\tau}{s}{\MPrin}}{p}$, where $\rulename{Let-Gen-BudgetError}$ converts an over-budget heal context to the fail-closed sink $\errV$ uniformly across all $\tau$. T4b/T4c cover Schema, arrow, and Policy types; the analogous $\mathtt{T4\_truthful\_success\_arrow}$ and $\mathtt{T4\_progress\_gen\_policy}$ are mechanized in \texttt{Safety.lean} as 1-line corollaries of the parametric $\mathtt{T4\_truthful\_success\_gen\_healable}$ / $\mathtt{T4\_progress\_gen\_healable}$.

Termination (Theorem in Appendix~\ref{app:term}) is argued on paper: between oracle firings only pure-core rules fire, which terminate under the loop assumption; at each oracle firing either $\MErr$ is flushed to $\emptyset$ (success) or an $\ErrEv{\cdot}$ extends it by one (heal), bounding heals by $\MRetry$ per site. $\rulename{Let-Gen-BudgetError}$ terminates the run at the first over-budget configuration by stepping to $\errV$; $\rulename{Let-Gen-TypeError}$ does the same at non-healable $\tau$. At a healable $\tau$ an ill-typed oracle response has no transition --- the trace is stuck, and Progress (T4c) rules this out under eventually-truthful $\mathcal{O}$.

\noindent\textbf{Scope of the formal results.}
T1--T4c are mechanized in Lean~4 (\texttt{mechanization/HADL/}), sorry-free and depending only on the standard meta-axioms $\mathtt{propext}$, $\mathtt{Classical.choice}$, $\mathtt{Quot.sound}$; a build-time guard (\texttt{AxiomCheck.lean}) breaks the build on any further axiom. The trace-level termination argument is prose-only (its per-step ingredients are mechanized). The eventually-truthful oracle is an assumption about LLM behaviour the runtime cannot verify; progress conclusions are conditional on it, and absent it a configuration may be stuck at an oracle redex whose budget is exhausted --- by design, \hadl fails closed by absence of transition. Static type checking is black-boxed (only reflexivity cases appear in proofs); the Rust type checker is not part of the Lean artifact. The threat-model gaps of \sect{sec:design-threat} --- prompt injection, semantic policy misalignment, \texttt{say}-based exfiltration, and \texttt{js} side effects --- are outside the scope of soundness. Practical claims beyond staged type safety and policy monotonicity per step are the subject of \sect{sec:evaluation}.
